\clearpage
\phantomsection
\section*{ЗАКЛЮЧЕНИЕ}
\addcontentsline{toc}{section}{ЗАКЛЮЧЕНИЕ}
В заключении можно выделить следующие результаты работы:
\begin{enumerate}
	\item В ходе анализа предметной области были описаны сферы ресторанной деятельности и предложены возможности оптимизации повседневных задач. Кроме того, были выявлены достоинства и недостатки существующих систем управления ресторанным бизнесом, а также установлена необходимость использования программно-информационных систем в сфере заведений общественного питания. 
	\item Составление диаграммы компонентов помогло визуализировать архитектуру проектируемой системы управления и установить связи между компонентами, что привело к ускорению внесения изменений и упрощению тестирование продукта.
	\item Разработка модели данных проектируемой системы способствовала организации и структурированию  информации, необходимой для проектирования базы данных.
	\item Создание диаграммы взаимодействий пользователя с системой позволила описать функциональные требования создаваемого продукта. Также были определены требования к входным и выходным данным.  
	\item Было произведено проектирование программно-информационной системы ресторанного бизнеса, разработка архитектуры  и создание интуитивно понятного пользовательского интерфейса программно-информационной системы управления ресторанным бизнесом.
	\item Проведено системное тестирование системы.
\end{enumerate}

Разработанная система представляет собой комплексное решение, состоящее из модулей для автоматизации различных направлений ресторанной деятельность. Также были реализованы функции создания, учёта и обработки заказов, анализ финансовых показателей, учёт данных о сотрудниках и позициях меню. Благодаря интеграции всех вышеперечисленных компонентов обеспечивается единое программно-информационное пространство, способствующее снижению ошибок и понижению времени обработки заказов.

Все требования, объявленные в техническом задании, были полностью реализованы, все задачи, поставленные в начале разработки проекта, были решены в полном объеме.


Разработанный программный продукт может быть использован для автоматизации ресторанного бизнеса или как основа для создания коммерческих программных продуктов.

Перспективой дальнейшей разработки является расширение программного функционала, создание дополнительных web-приложений и мобильных программ.

